% =====================================================================
\section{Interrupts, exceptions and reset}
\label{sec:int}
% =====================================================================

The processor has three non-maskable exceptions (WOVF, WUNF, ILL) and two
maskable interrupts (INT0, \IIC), served through a fixed vector table at the
start of program memory; the user program starts at \hex{0020}.

\subsection{Vector table}

\begin{center}
\small
\begin{tabularx}{\linewidth}{@{}l X l l@{}}
  \toprule
  \textbf{Address} & \textbf{Source} & \textbf{Type} & \textbf{Priority} \\
  \midrule
  \hex{0000} & RESET & -- & highest \\
  \hex{0004} & WOVF: register window overflow & exception, non-maskable & 1 \\
  \hex{0008} & WUNF: register window underflow & exception, non-maskable & 2 \\
  \hex{000C} & ILL: illegal instruction / \texttt{TRAP} & exception, non-maskable & 3 \\
  \hex{0010} & INT0: external interrupt & maskable & 4 \\
  \hex{0014} & \IIC: operation complete & maskable & 5 \\
  \hex{0018}--\hex{001F} & Reserved & -- & -- \\
  \bottomrule
\end{tabularx}
\end{center}

Each vector is 4 bytes long, which fits a \texttt{JMP addr16} instruction (3 bytes).

\subsection{Requirements}

\begin{reqtable}
\req{INT-001}{A maskable interrupt \MUST{} be taken only when SR.IE = 1, its
  bit in INTEN is set and its flag in INTFLG is set.}
\req{INT-002}{Interrupts \MUST{} be taken only between instructions;
  exceptions (WOVF, WUNF, ILL) are raised before execution of the instruction
  that caused them.}
\req{INT-003}{On entry to a handler, the processor \MUST{} push PC (2 bytes)
  and SR (1 byte) onto the stack, clear SR.IE and jump to the vector address.}
\req{INT-004}{For an interrupt, the saved PC \MUST{} point to the next
  instruction; for an exception, to the instruction that raised it.}
\req{INT-005}{\texttt{RETI} \MUST{} pop SR and PC from the stack (restoring IE).}
\req{INT-006}{When several sources are pending, the source with the highest
  priority according to the vector table \MUST{} be served first.}
\req{INT-007}{The latency from the end of the current instruction to the first
  instruction of the handler \SHOULD{} not exceed 8 cycles; the actual value
  \MUST{} be documented.}
\req{INT-008}{The INT0 input \MUST{} be synchronised to CLK; the interrupt is
  triggered on the falling edge.}
\req{INT-009}{A handler that uses \texttt{WINC} \MUST{} execute \texttt{WDEC}
  before \texttt{RETI}; the WOVF/WUNF handlers \MUSTNOT{} use \texttt{WINC} or
  \texttt{WDEC}.}
\req{INT-010}{\texttt{HALT} \MUST{} suspend execution until an interrupt occurs
  (also with IE = 0, in which case execution resumes at the next instruction
  without entering the handler).}
\req{INT-011}{After reset: PC = \hex{0000}, SR = \hex{00}, CWP = WSTART = 0,
  INTEN = INTFLG = 0, \IIC{} controller disabled, SP = \hex{0800} (the
  first \texttt{PUSH} writes the last byte of data memory, \hex{07FF}).}
\req{INT-012}{The \texttt{RST\_N} signal \MUST{} be active low, with
  asynchronous assertion and synchronous de-assertion (2 flip-flops).}
\end{reqtable}
